\section{正合三角}\label{sec:3}
\subsection{基本思路}\label{sec:3a}
在任何一种范畴中构造对象时，常见的策略都是：先用所期望的表示函子刻画所需对象，再研究存在性（可表示性）问题。在我们的情形中，适当的形式体系是 $\Ai$ 模与 Yoneda 嵌入。固定一个 $\Ai$ 范畴 $\A$，记 $\Q=\modcat(\A)$。设 $\M$ 是一个 $\Ai$ 模。假设存在 $Y\in\Ob\A$，其 Yoneda 嵌入下的像 $\Y$ 在 $H^0(\Q)$ 中与 $\M$ 同构，并固定该范畴中的一个同构 $[t]:\Y\to\M$。
\begin{definition*}
称二元组 $(Y,[t])$ \emph{拟表示} $\M$。
\end{definition*}
\stref{推论}{2.13}表明，若 $(\widetilde Y,[\tilde t])$ 是另一个这样的二元组，那么在 $H^0(\A)$ 中存在一个优选的同构 $[a]:Y\to\widetilde Y$，满足 $[\tilde t]\cdot[\mathcal I^1(a)]=[t]$。\znote{32}{“拟表示”把通常的表示概念在同伦意义下放宽：并不要求 $\hom_{\Q}(\Y,\M)$ 与 $\M(Y)$ 同构，只要求 Yoneda 像 $\Y$ 与 $\M$ 在 $H^0(\Q)$ 中同构。这样在作映射锥、平移等构造时，每个对象只需在“相差一个同构”的意义下被表示，正是构造三角化包络所需要的灵活性。[来源：译者按（上下文推断，无外部依据）：对照\stref{引理}{3.1}与\stref{推论}{2.13}]}略微滥用术语，有时我们会直接称 $Y$ 为拟表示 $\M$ 的对象，但仍把 $[t]$ 的选择视为数据的一部分。利用\eref{1.25}中的链映射 $\lambda=\lambda_{\M}:\M(Y)\to\hom_{\Q}(\Y,\M)$，可以给出更具体的刻画：
\begin{lemma}\label{st:3.1}
$(Y,[t])$ 拟表示 $\M$，当且仅当如下条件成立：存在 $c\in\M^0(Y)$，满足 $\mu^1_{\M}(c)=0$ 以及 $[t]=[\lambda(c)]$，并且对所有 $X\in\Ob\A$，映射
\begin{equation}\label{eq:3.1}
\hom_{\A}(X,Y)\longrightarrow\M(X),\qquad b\longmapsto(-1)^{|b|}\mu^2_{\M}(c,b)
\end{equation}
都是拟同构。
\end{lemma}
\begin{proof}
由于 $\mu^1_{\M}(c)=0$，$\lambda(c)$ 是一个 $\Ai$ 模同态 $\Y\to\M$，其一阶项除符号外就是\eref{3.1}。\stref{引理}{1.16}表明 $[t]=[\lambda(c)]$ 在 $H^0(\Q)$ 中是同构，因此 $(Y,[t])$ 拟表示 $\M$。反过来，若已有拟表示二元组 $(Y,[t])$，那么由\stref{引理}{2.12}，存在 $c$ 满足 $\mu^1_{\M}(c)=0$ 及 $[t]=[\lambda(c)]$。由于 $[\lambda(c)]$ 在 $H^0(\Q)$ 中是同构，它必然在基础上同调模之间诱导同构，这就意味着\eref{3.1}对所有 $X$ 都是拟同构。
\end{proof}

\subsection{直和}\label{sec:3b}
下面是这种方法的一个初等例子。给定 $\Ai$ 模 $\M_0,\M_1$，把相应的分次向量空间取直和 $\M_0(X)\oplus\M_1(X)$（$X\in\Ob\A$），并赋予结构映射 $\mu^d_{\M_0}\oplus\mu^d_{\M_1}$，即可定义它们的直和 $\M_0\oplus\M_1$。现在，设 $Y_0,Y_1$ 为 $\A$ 的对象，$\Y_0,\Y_1$
\SourcePage{41}
为其 Yoneda 嵌入下的像。
\begin{definition*}
把任意一个拟表示 $\Ai$ 模 $\Y_0\oplus\Y_1$ 的对象记为 $Y_0\oplus Y_1$。
\end{definition*}
由定义，对每个 $X$，它都带有典范同构
\[
\Hom_{H(\A)}(X,Y_0\oplus Y_1)\cong\Hom_{H(\A)}(X,Y_0)\oplus\Hom_{H(\A)}(X,Y_1).
\]

\subsection{张量积}\label{sec:3c}
同样在这个基本层次上，设 $Z=(Z,\partial_Z)$ 是 $\K$ 向量空间的链复形，$\M$ 是 $\Ai$ 模。
\begin{definition*}
张量积模 $Z\otimes\M$ 定义为
\begin{equation}\label{eq:3.2}
\begin{aligned}
(Z\otimes\M)(X)&=Z\otimes\M(X),\\
\mu^1_{Z\otimes\M}(z\otimes b)&=(-1)^{|b|-1}\partial_Z(z)\otimes b+z\otimes\mu^1_{\M}(b),\\
\mu^d_{Z\otimes\M}(z\otimes b,a_{d-1},\ldots,a_1)&=z\otimes\mu^d_{\M}(b,a_{d-1},\ldots,a_1)\quad(d\geq2).
\end{aligned}
\end{equation}

\end{definition*}
可以使 $Z\otimes-$ 成为 $\Q$ 到自身的 $\Ai$ 函子。在一阶上，它把 $t\in\hom_{\Q}(\M_0,\M_1)$ 映到 $\id_Z\otimes t\in\hom_{\Q}(Z\otimes\M_0,Z\otimes\M_1)$，后者由下式定义：
\[
(\id_Z\otimes t)^d(z\otimes b,a_{d-1},\ldots,a_1)=(-1)^{|z|(|t|-1)}z\otimes t^d(b,a_{d-1},\ldots,a_1);
\]
所有高阶项均置为零。给定 $\A$ 的对象 $Y$，其 Yoneda 像为 $\Y$；把任意一个拟表示 $Z\otimes\Y$ 的对象记为 $Z\otimes Y$。在上同调层次上，有
\[
\begin{aligned}
\Hom_{H(\A)}(X,Z\otimes Y)&\cong\Hom_{H(\Q)}(X,Z\otimes\Y)\cong H((Z\otimes\Y)(X))\\
&\cong H(Z)\otimes H(\Y(X))\cong H(Z)\otimes\Hom_{H(\A)}(X,Y),
\end{aligned}
\]
其中的同构分别来自 Yoneda 嵌入、映射 $\lambda$ 以及通常的 Künneth 公式。

\begin{note}\label{st:3.2}
给定链复形 $(Z,\partial_Z)$，令 $\widetilde Z=H(Z)$ 且 $\partial_{\widetilde Z}=0$，再选择一个线性映射 $\psi:\widetilde Z\to Z$，为每个上同调类选取链代表。由\stref{引理}{1.16}，所得嵌入 $\widetilde Z\otimes\M\hookrightarrow Z\otimes\M$ 在 $H^0(\Q)$ 中诱导同构（而且这个同构不依赖于 $\psi$ 的选择）。这说明，把张量积运算限制在 $Z$ 只是一个微分为零的分次 $\K$ 向量空间的情形，并无本质上的一般性损失。
\end{note}

\subsection{平移}\label{sec:3d}
\begin{definition*}
最常用的张量积例子，是 $Z$ 为集中在某个次数 $-\sigma$ 的一维向量空间 $\K[\sigma]$ 的情形。所得对象分别记为 $\K[\sigma]\otimes\M=\ModShift^{\sigma}\M$ 与 $\K[\sigma]\otimes Y=S^{\sigma} Y$，称为 $\M$ 或 $Y$ 的
\SourcePage{42}
$\sigma$ 次\emph{平移}。当 $\sigma=1$ 时省略指标，只写 $\ModShift\M$、$SY$。
\end{definition*}
这些平移如果存在，就带有典范同构
\begin{align}
\Hom_{H(\A)}(Y_0,S^{\sigma} Y_1)&\cong\Hom_{H(\A)}(Y_0,Y_1)[\sigma],\label{eq:3.3}\\
\Hom_{H(\A)}(Y_0,Y_1)&\cong\Hom_{H(\A)}(S^{\sigma} Y_0,S^{\sigma} Y_1),\label{eq:3.4}\\
\Hom_{H(\A)}(S^{\sigma} Y_0,Y_1)[\sigma]&\cong\Hom_{H(\A)}(Y_0,Y_1).\label{eq:3.5}
\end{align}
第一个来自定义，第二个来自 $\ModShift^{\sigma}=\K[\sigma]\otimes-$ 的函子性。第三个则先对二元组 $(S^{\sigma} Y_0,Y_1)$ 应用\eref{3.3}，再对 $(Y_0,Y_1)$ 反向（从右到左）应用\eref{3.4}而得。\znote{34}{这类符号差异在三角范畴中是常态：一个三角 $(f,g,h)$ 换成 $(f,g,-h)$ 一般不再是殊三角，尽管它仍诱导出长正合列；而 $(f,-g,-h)$ 因与原三角同构而仍是殊三角。因此使用平移同构时必须固定符号约定，否则不同来源的公式会相差 $(-1)^{\sigma}$。[来源：nLab "triangulated category"，https:/\allowbreak /\allowbreak ncatlab.org/\allowbreak nlab/\allowbreak show/\allowbreak triangulated+\allowbreak category（2026-11-15 访问）]}

\clearpage
\begin{note}\label{st:3.3}
得到\eref{3.5}的另一种途径，是从 $(S^{\sigma} Y_0,Y_1)$ 出发，应用\eref{3.4}得到 $(Y_0,S^{-\sigma}Y_1)$，然后反向使用\eref{3.3}回到 $(Y_0,Y_1)$。遗憾的是，这与前面的定义相差一个符号 $(-1)^{\sigma}$，因此使用这些同构时需要稍加留意。
\end{note}

假设 $\Ai$ 范畴 $\A$ 满足：对所有 $Y$，$SY$ 都存在。那么可以构造一个平移 $\Ai$ 函子 $S:\A\to\A$，它在 $H^0(\operatorname{fun}(\A,\A))$ 中的同构意义下唯一，并使图
\SourceDiagram{diagram-3d}{0.13915370}
在 $H^0(\operatorname{fun}(\A,\Q))$ 中在同构意义下交换。具体地，考虑由所有形如 $SY$ 的对象组成的全 $\Ai$ 子范畴 $\widetilde{\A}\subset\A$；相应地，令 $\widetilde{\Q}\subset\Q$ 为由如下 $\Ai$ 模组成的子范畴：它们在 $H^0(\Q)$ 中同构于某个 $\ModShift\Y=\ModShift\mathcal I(Y)$。由构造，Yoneda 嵌入的限制是拟等价 $\widetilde{\A}\to\widetilde{\Q}$，故由\stref{定理}{2.9}可取其逆。给定一个这样的逆 $\mathcal K$，令 $S=\mathcal K\circ\ModShift\circ\mathcal I$ 即可。这里以平移为例作了说明，但同样的论证适用于 $\Ai$ 模的其他函子性构造。\znote{34}{这里构造的 $S$ 就是经典三角范畴中的平移函子：代数情形常记 $(-)[1]$，也写作 $S$ 或 $\Sigma$，它在三角范畴中应当是范畴的自等价。本节的论证说明，只要每个对象的平移都存在，就能把这一运算提升为 $\Ai$ 函子，从而在 $H^0$ 上得到真正的自等价。[来源：nLab "triangulated category"，https:/\allowbreak /\allowbreak ncatlab.org/\allowbreak nlab/\allowbreak show/\allowbreak triangulated+\allowbreak category（2026-11-15 访问）]}

\begin{note}\label{st:3.4}
任意平移 $S^{\sigma} Y$（$\sigma\in\Z$）和直和 $Y_0\oplus Y_1$ 的存在性，蕴含对于上同调有限维的复形 $(Z,\partial_Z)$，张量积 $Z\otimes Y$ 也存在。具体地，在 $\widetilde Z=H(Z)$ 中选取一组基 $\{[z^i]\}$，令 $\sigma^i=-|z^i|$，则
\[
\bigoplus_{i\in I}S^{\sigma^i}Y
\]
在 Yoneda 嵌入下映到 $\bigoplus_{i\in I}\ModShift^{\sigma^i}\Y\cong\widetilde Z\otimes\Y$。由\stref{笔记}{3.2}，它也是 $Z\otimes\Y$ 的拟表示对象。
\end{note}
\SourcePage{43}
\subsection{锥}\label{sec:3e}
设 $Y_0,Y_1$ 为 $\A$ 的对象，$c\in\hom^0_{\A}(Y_0,Y_1)$ 为零次上闭链（$\mu^1_{\A}(c)=0$）。
\begin{definition*}
$c$ 的\emph{抽象映射锥}是如下定义的 $\Ai$ 模 $\C=\ACone(c)$：
\begin{equation}\label{eq:3.6}
\begin{aligned}
\C(X)&=\hom_{\A}(X,Y_0)[1]\oplus\hom_{\A}(X,Y_1),\\
\mu^d_{\C}((b_0,b_1),a_{d-1},\ldots,a_1)
&=\bigl(\mu^d_{\A}(b_0,a_{d-1},\ldots,a_1),\\[-2pt]
&\qquad\mu^d_{\A}(b_1,a_{d-1},\ldots,a_1)+\mu^{d+1}_{\A}(c,b_0,a_{d-1},\ldots,a_1)\bigr).
\end{aligned}
\end{equation}

\end{definition*}
直接计算即可验证\eref{1.19}。至于 $\C$ 的 c-含幺性，需要检验：若 $e_X\in\hom^0_{\A}(X,X)$ 代表 $H(\A)$ 中的恒等态射，则 $(b_0,b_1)\mapsto\mu^2_{\C}((b_0,b_1),e_X)$ 在 $H(\C(X),\mu^1_{\C})$ 上诱导的自同态是恒等映射。我们已经知道它是幂等的；考察 $\C(X)$ 中由 $b_0=0$ 给出的子复形及其商，又可知它可逆，因此它就是恒等映射。
\begin{definition*}
把拟表示抽象映射锥的对象记为 $\Cone(c)$，简称 $c$ 的\emph{映射锥}。
\end{definition*}
与通常的情形一样，这样的对象如果存在，就在 $H^0(\A)$ 中的典范同构意义下唯一。\znote{35}{抽象映射锥以链复形的映射锥为原型：对链映射 $f:X\to Y$，$\Cone(f)$ 的第 $n$ 项是 $Y_n\oplus X_{n-1}$，其微分的左上、右下两块分别是两个复形的微分，右上块为 $f$，左下块为零，且右下块带负号。把 $X$ 换成 $\hom_{\A}(X,Y_0)$、$f$ 换成与 $c$ 复合，即得\eref{3.6}，其中 $c$ 前的符号也源于这一约定。[来源：nLab "mapping cone"，https:/\allowbreak /\allowbreak ncatlab.org/\allowbreak nlab/\allowbreak show/\allowbreak mapping+\allowbreak cone（2026-11-15 访问）]}

考察在一个固定上同调类内改变上闭链 $c$ 的选择时，$\ACone(c)$ 如何变化，是很有启发性的。若另取代表 $\tilde c=c+\mu^1_{\A}(h)$，则
\begin{equation}\label{eq:3.7}
\begin{aligned}
t^1(b_0,b_1)&=\bigl((-1)^{|b_0|-1}b_0,\,(-1)^{|b_1|}b_1+\mu^2_{\A}(h,b_0)\bigr),\\
t^d((b_0,b_1),a_{d-1},\ldots,a_1)&=\bigl(0,\mu^{d+1}_{\A}(h,b_0,a_{d-1},\ldots,a_1)\bigr)\quad(d\geq2)
\end{aligned}
\end{equation}
定义了模同态 $\ACone(c)\to\ACone(\tilde c)$。由\stref{引理}{1.16}，它在 $H^0(\Q)$ 中是同构（这里及下文类似情形中，$|b_0|$ 指 $b_0$ 在 $\hom_{\A}(X,Y_0)$ 中的次数，即施加平移 $[1]$ 之前的次数）。因此，上同调类 $[c]$ 在该范畴中的非典范同构意义下确定 $\ACone(c)$；之所以非典范，是因为\eref{3.7}需要选择 $h$。由定义，对作为 $H^0(\A)$ 对象的 $\Cone(c)$ 也成立同样的结论。

\begin{note}\label{st:3.5}
直和、平移与映射锥之间有一个初等关系：平凡态射 $c=0\in\hom_{\A}(Y_0,Y_1)$ 的锥，拟表示 $\ModShift\Y_0\oplus\Y_1$。例如，假设 $\A$ 对映射锥封闭，也就是说，对所有 $c$，$\Cone(c)$ 都存在。先取一个辅助对象 $Z$，以及代表恒等态射的链 $e_Z$，则 $\ACone(e_Z)$ 同构于零模。于是，对任意 $Y$，零映射 $Y\to\Cone(e_Z)$ 的锥其实就是 $SY$，这意味着 $\A$ 对平移 $S^{\sigma}$（$\sigma>0$）封闭（显然也可推出，对所有 $Y_0,Y_1$，$SY_0\oplus SY_1$ 都存在）。
\end{note}

\subsection{正合三角}\label{sec:3f}
任意抽象映射锥 $\C=\ACone(c)$ 都带有典范模同态 $\iota\in\hom^0_{\Q}(\Y_1,\C)$、$\pi\in\hom^1_{\Q}(\C,\Y_0)$；用\eref{3.6}中的记号，
\SourcePage{44}
它们由 $\iota^1(b_1)=(0,(-1)^{|b_1|}b_1)$、$\pi^1(b_0,b_1)=(-1)^{|b_0|-1}b_0$ 给出，高阶项为零。这些态射在 $H(\Q)$ 中组成三角图：
\SourceEquationDiagram{diagram-3-8}{0.21807669}{3.8}
箭头上的标记 $[1]$ 提醒我们，相应态射的次数为一。为便于后文引用，我们还要给出这些映射的另一些代表：
\begin{lemma}\label{st:3.6}
设 $e_{Y_0},e_{Y_1}$ 是代表 $H(\A)$ 中恒等态射的链。那么由
\[
\begin{aligned}
\tilde\iota^d(b,a_{d-1},\ldots,a_1)&=\bigl(0,\mu^{d+1}_{\A}(e_{Y_1},b,a_{d-1},\ldots,a_1)\bigr),\\
\tilde\pi^d((b_0,b_1),a_{d-1},\ldots,a_1)&=-\mu^{d+1}_{\A}(e_{Y_0},b_0,a_{d-1},\ldots,a_1)
\end{aligned}
\]
给出的模同态 $\tilde\iota\in\hom^0_{\Q}(\Y_1,\C)$、$\tilde\pi\in\hom^1_{\Q}(\C,\Y_0)$，分别与 $\iota$、$\pi$ 上同调。
\end{lemma}
\begin{proof}
已知 Yoneda 嵌入上同调含幺。因此存在 $h\in\hom^{-1}_{\Q}(\Y_1,\Y_1)$，使 $\mu^1_{\Q}(h)=e_{\Y_1}-\mathcal I^1(e_{Y_1})$，即 $\Q$ 中的严格恒等态射与来自 $\A$ 中恒等态射代表的链之差。把它作为第二个分量，第一个分量置零，就得到 $k=(0,h)\in\hom^{-1}_{\Q}(\Y_1,\C)$，显然满足 $\mu^1_{\Q}(k)=\iota-\tilde\iota$。对 $\pi-\tilde\pi$ 的论证相同。
\end{proof}

\begin{definition*}
把 $H(\A)$ 中任意如下形式的图
\SourceEquationDiagram{diagram-3-9}{0.21807669}{3.9}
称为\emph{正合三角}，如果它在 Yoneda 嵌入下变为与\eref{3.8}同构的图。
\end{definition*}
这意味着：令 $c=c_1$，则在 $H^0(\Q)$ 中存在同构 $[t]:\Y_2\to\C=\ACone(c)$，满足 $[\pi]\circ[t]=[\mathcal I^1(c_3)]$ 以及 $[\iota]=[t]\circ[\mathcal I^1(c_2)]$。根据前面关于 $\ACone(c)$ 如何依赖链 $c$ 的选择的讨论，这里的选择并不重要：若对某个 $c$ 的选择能构造同构 $[t]$，那么对任何其他选择也都能做到。沿用\stref{引理}{3.1}的一般思路，可以得到更具体的正合性判据：

\SourcePage{45}
\begin{lemma}\label{st:3.7}
三角\eref{3.9}正合，当且仅当如下条件成立。存在
\[
h_1\in\hom^0_{\A}(Y_1,Y_0),\qquad h_2\in\hom^0_{\A}(Y_2,Y_1),\qquad k\in\hom^{-1}_{\A}(Y_1,Y_1),
\]
满足
\begin{equation}\label{eq:3.10}
\begin{aligned}
\mu^1_{\A}(h_1)&=\mu^2_{\A}(c_3,c_2),\\
\mu^1_{\A}(h_2)&=-\mu^2_{\A}(c_1,c_3),\\
\mu^1_{\A}(k)&=-\mu^2_{\A}(c_1,h_1)+\mu^2_{\A}(h_2,c_2)+\mu^3_{\A}(c_1,c_3,c_2)-e_{Y_1},
\end{aligned}
\end{equation}
其中 $e_{Y_1}$ 是恒等态射的链代表；此外，对每个对象 $X$，如下链复形必须无上同调：
\begin{equation}\label{eq:3.11}
\begin{gathered}
\hom_{\A}(X,Y_2)[1]\oplus\hom_{\A}(X,Y_0)[1]\oplus\hom_{\A}(X,Y_1),\\
\partial=\begin{pmatrix}
\mu^1_{\A}&0&0\\
\mu^2_{\A}(c_3,\cdot)&\mu^1_{\A}&0\\
\mu^2_{\A}(h_2,\cdot)+\mu^3_{\A}(c_1,c_3,\cdot)&\mu^2_{\A}(c_1,\cdot)&\mu^1_{\A}
\end{pmatrix}.
\end{gathered}
\end{equation}
\end{lemma}
\begin{proof}
假设上述所有数据都存在。令 $\C=\ACone(c_1)$，则 $b=(c_3,h_2)$ 是 $\C(Y_2)$ 中的上闭链，从而给出模同态 $t=\lambda_{\C}(b)\in\hom^0_{\Q}(\Y_2,\C)$。复形\eref{3.11}就是相关链映射\eref{3.1}的映射锥，因此其无上同调性保证 $[t]$ 在 $H^0(\Q)$ 中是同构。由\eref{1.26}及 $\C$ 的定义可知，$[\mu^2_{\Q}(t,\mathcal I^1(c_2))]=[\lambda_{\C}(\mu^2_{\C}(b,c_2))]$ 是
\[
\bigl(\mu^2_{\A}(c_3,c_2),\,\mu^2_{\A}(h_2,c_2)+\mu^3_{\A}(c_1,c_3,c_2)\bigr)\in\C(Y_1)
\]
在 $\lambda_{\C}$ 下的像。由假设\eref{3.10}，它与 $(0,e_{Y_1})$ 上同调，而后者在 $\lambda_{\C}$ 下的像，就是\stref{引理}{3.6}中的模同态 $\tilde\iota$。因此
\[
[\iota]=[\mu^2_{\Q}(t,\mathcal I^1(c_2))].
\]
接着，由\eref{1.27}可见，在 $\Hom^1_{H(\Q)}(\Y_2,\Y_0)$ 中，
\[
[\mu^2_{\Q}(\pi,t)]=[\lambda_{\Y_0}(\pi^1(b))]=[\lambda_{\Y_0}(c_3)]=[\mathcal I^1(c_3)].
\]
这样就验证了正合三角所需的全部条件。\znote{37}{引理 3.7 的条件正对应经典三角范畴的公理：\eref{3.10}的前两式记录“三角中相继复合为零”，而\eref{3.11}无上同调则等价于“用 Hom 作用后得到长正合列”。后者是三角结构最核心的性质，在通常的公理体系中单列为一条引理。[来源：nLab "triangulated category"，https:/\allowbreak /\allowbreak ncatlab.org/\allowbreak nlab/\allowbreak show/\allowbreak triangulated+\allowbreak category（2026-11-15 访问）]}

反方向的证明，实质上就是把上述论证倒过来。任意态射 $[t]:\Y_2\to\C$ 都经由 $\lambda_{\C}$ 来自某个元素 $b\in\C(Y_2)$。由于 $[\mathcal I^1(\pi^1(b))]=[\mu^2_{\Q}(\pi,t)]=[\mathcal I^1(c_3)]$，$b$ 的第一个分量与 $c_3$ 上同调。加上一个适当的上边界后，可以假设 $b=(c_3,h_2)$，其中 $h_2$ 为某个元素。接着，
\[
\begin{aligned}
{}[\lambda_{\C}(0,e_{Y_1})]&=[\tilde\iota]=[\mu^2_{\Q}(t,\mathcal I^1(c_2))]\\
&=[\lambda_{\C}(\mu^2_{\C}(b,c_2))].
\end{aligned}
\]
因为 $\lambda_{\C}$ 是拟同构，$(0,e_{Y_1})$ 与 $\bigl(\mu^2_{\A}(c_3,c_2),\mu^2_{\A}(h_2,c_2)+\mu^3_{\A}(c_1,c_3,c_2)\bigr)$ 上同调。选择一个余链，其微分为这两个上闭链之差，即得到剩余数据 $(k,h_1)$。
\end{proof}

\SourcePage{46}
\subsection{另一种刻画}\label{sec:3g}
虽然拟表示对象这一概念为各种构造提供了统一的方法，但用它推导具体性质时，可能会显得繁琐。因此，接下来关于正合三角的讨论暂时不给出证明。稍后改用更具体的观点时，我们会补上这些证明。

考虑对象为 $Z_0,Z_1,Z_2$ 的严格含幺 $\Ai$ 范畴 $\mathcal D$，其中
\[
\begin{aligned}
\hom_{\mathcal D}(Z_k,Z_k)&=\K\cdot e_{Z_k},\\
\hom_{\mathcal D}(Z_0,Z_1)&=\K\cdot x_1,\quad |x_1|=0,\\
\hom_{\mathcal D}(Z_1,Z_2)&=\K\cdot x_2,\quad |x_2|=0,\\
\hom_{\mathcal D}(Z_2,Z_0)&=\K\cdot x_3,\quad |x_3|=1,\\
\hom_{\mathcal D}(Z_2,Z_1)&=\hom_{\mathcal D}(Z_1,Z_0)=\hom_{\mathcal D}(Z_0,Z_2)=0;
\end{aligned}
\]
除严格单位元所要求的运算外，仅有的非零复合映射为
\[
\mu^3_{\mathcal D}(x_3,x_2,x_1)=e_{Z_0},\quad
\mu^3_{\mathcal D}(x_1,x_3,x_2)=e_{Z_1},\quad
\mu^3_{\mathcal D}(x_2,x_1,x_3)=e_{Z_2}.
\]
\begin{proposition}\label{st:3.8}
$H(\A)$ 中的三角\eref{3.9}正合，当且仅当存在 $\Ai$ 函子 $\F:\mathcal D\to\A$，使 $\F(Z_k)=Y_k$，并在上同调上满足 $[\F^1(x_k)]=[c_k]$。
\end{proposition}
将它与\stref{引理}{3.7}比较很有启发性。假设存在如上所述的 $\Ai$ 函子 $\F$。选择 $d_k$，使 $\mu^1_{\A}(d_k)=\F^1(x_k)-c_k$。那么
\begin{equation}\label{eq:3.12}
\begin{aligned}
h_1&=-\F^2(x_3,x_2)+\mu^2_{\A}(\F^1(x_3),d_2)-\mu^2_{\A}(d_3,c_2),\\
h_2&=\F^2(x_1,x_3)-\mu^2_{\A}(\F^1(x_1),d_3)-\mu^2_{\A}(d_1,c_3),\\
k&=-\mu^3_{\A}(d_1,c_3,c_2)-\mu^3_{\A}(\F^1(x_1),d_3,c_2)+\mu^3_{\A}(\F^1(x_1),\F^1(x_3),d_2)\\
&\quad-\mu^2_{\A}(\F^2(x_1,x_3),d_2)\\
&\quad+\mu^2_{\A}\bigl(d_1,\mu^2_{\A}(\F^1(x_3),d_2)-\mu^2_{\A}(d_3,c_2)-\F^2(x_3,x_2)\bigr)
\end{aligned}
\end{equation}
满足\eref{3.10}中的全部等式，其中 $e_{Y_1}=\F^1(e_{Z_1})$。$\F$ 的其他分量可用来构造复形\eref{3.11}的收缩同伦（这里不把它写出），这给出了证明\stref{命题}{3.8}中一个方向的一种方法。稍后给出的证明在原则上与之相似，但会有更好的理论基础，从而减少对\eref{3.12}这类显式公式的依赖。
\begin{corollary}\label{st:3.9}
若 $\F:\A\to\B$ 为任意 $\Ai$ 函子，则 $H(\F)$ 把 $H(\A)$ 中的正合三角映为 $H(\B)$ 中的正合三角。
\end{corollary}
\begin{corollary}\label{st:3.10}
设 $\F:\A\to\B$ 为上同调全忠实 $\Ai$ 函子。假设 $H(\B)$ 中有一个正合三角，组成它的对象都落在 $\F$ 的像中。那么它的原像是 $H(\A)$ 中的正合三角。
\end{corollary}

\begin{note}\label{st:3.11}
\stref{命题}{3.8}出自 Kontsevich。我们的 $\mathcal D=\mathcal D_3$ 实际上属于一族 $\Ai$ 范畴 $\mathcal D_n$（$n\geq1$），后者有 $n$ 个对象 $Z_0,\ldots,Z_{n-1}$。
\SourcePage{47}
存在满足 $\F(Z_0)=Y_0$ 的 $\Ai$ 函子 $\F:\mathcal D_1\to\A$，当且仅当 $Y_0$ 是 $H(\A)$ 中的零对象。类似地，存在满足 $\F(Z_k)=Y_k$ 的 $\F:\mathcal D_2\to\A$，当且仅当 $Y_0,Y_1$ 在 $H^0(\A)$ 中同构（有关讨论参看\bcite{95}，第 193 页）。因此，正合三角可视为把同构概念从两个对象推广到三个对象。
\end{note}

\begin{note}\label{st:3.12}
还有一个值得注意的推论。取 $\mathcal D$ 中由对象 $Z_k$ 组成的基本正合三角。按\stref{笔记}{1.2}定义的任意三个相继映射的 Massey 积，除一个符号外就是恒等态射：
\begin{equation}\label{eq:3.13}
\begin{aligned}
\langle[x_3],[x_2],[x_1]\rangle&=[e_{Z_0}],\\
\langle[x_1],[x_3],[x_2]\rangle&=[-e_{Z_1}],\\
\langle[x_2],[x_1],[x_3]\rangle&=[e_{Z_2}].
\end{aligned}
\end{equation}
由 Massey 积对于 $\Ai$ 函子的函子性，任意正合三角都满足同样的结论。例如，取一个正合三角，把其中三个态射之一乘以某个 $\K\setminus\{0,1\}$ 中的常数。一般而言，所得三角将不满足\eref{3.13}，因而不再正合。这里说“一般而言”，是因为 Massey 积取值的商群\eref{1.4}可能为零；在这种情形下，修改后的三角仍可能正合。这在经典的三角范畴框架中也成立，例如参看\bcite{20}，推论 1.1.10(ii)。
\end{note}

\subsection{三角化的 \texorpdfstring{$A_\infty$}{A∞} 范畴}\label{sec:3h}
\begin{definition*}
若一个 $\Ai$ 范畴非空（$\Ob\A\neq\varnothing$），且满足下列条件，则称它为\emph{三角化的}：
\begin{itemize}
\item $H^0(\A)$ 中每个态射 $[c_1]$ 都能补全为正合三角\eref{3.9}；
\item 对每个对象 $Y$，存在 $\widetilde Y$，使 $S\widetilde Y\cong Y$ 在 $H^0(\A)$ 中成立。
\end{itemize}

\end{definition*}
这里需要略作解释。由\stref{笔记}{3.5}，第一个条件意味着每个对象 $Y$ 都有平移 $SY$，因此第二个条件才有意义。用更具函子性的语言来说，第一个条件允许构造平移函子 $S:\A\to\A$（参看\secref{3d}），第二个条件则断言它是拟等价。\znote{39}{三角化的两条要求恰好是经典公理 TR2 与平移可逆性的范畴化：第一条保证任意态射都能补全为殊三角，第二条保证平移在两个方向上都有定义。对 $H^0(\A)$ 而言平移是自等价，所以第二条自动成立，真正的内容全在第一条。[来源：nLab "triangulated category"，https:/\allowbreak /\allowbreak ncatlab.org/\allowbreak nlab/\allowbreak show/\allowbreak triangulated+\allowbreak category（2026-11-15 访问）]}
\begin{example}\label{st:3.13}
$\Ai$ 范畴 $\Q=\modcat(\A)$ 是三角化的。
\end{example}
\begin{proposition}\label{st:3.14}
若 $\A$ 是三角化的 $\Ai$ 范畴，则 $H^0(\A)$ 是经典意义下的三角范畴。此外，若 $\A,\B$ 是三角化的 $\Ai$ 范畴，$\F:\A\to\B$ 是 $\Ai$ 函子，则 $H^0(\F)$ 是正合函子。
\end{proposition}
回顾经典意义下的三角范畴\bcite{141}、\bcite{61}：它是一个加性范畴，配以一个“平移”自等价 $T$ 以及一类“殊三角”，后者是如下形式的周期图：
\[
\cdots\longrightarrow Y_0\longrightarrow Y_1\longrightarrow Y_2\longrightarrow TY_0\longrightarrow TY_1\longrightarrow TY_2\longrightarrow\cdots.
\]
\SourcePage{48}
在我们的情形中，这个平移由平移函子诱导，即 $T=H^0(S)$；殊三角就是正合三角，其中借助\eref{3.3}把 $[c_3]$ 转化为态射 $Y_2\to SY_0$。\znote{40}{经典定义要求加性范畴、平移自等价与一族满足 (TR1)–(TR4) 的三角，其中 (TR4) 即八面体公理。本书取 $T=H^0(S)$，并把正合三角规定为殊三角，于是 $\Ai$ 层面的正合性与经典三角结构被衔接起来，后面各条推论都可以回到这套公理去检验。[来源：nLab "triangulated category"，https:/\allowbreak /\allowbreak ncatlab.org/\allowbreak nlab/\allowbreak show/\allowbreak triangulated+\allowbreak category（2026-11-15 访问）]}

\clearpage
\subsection{三角的性质}\label{sec:3i}
\stref{命题}{3.14}表明，三角化 $\Ai$ 范畴中的正合三角类满足三角结构的公理性质。事实上，这些性质经过适当表述后，也适用于一般的 $\Ai$ 范畴。为了方便不太熟悉同调代数的读者，下面列出其中一部分：
\begin{corollary}\label{st:3.15}
设\eref{3.9}是 $\Ai$ 范畴 $\A$ 中的正合三角。那么其中任意两个相继态射的复合为零；而且，固定对象 $X$，从左侧或右侧取 Hom，可得到向量空间的长正合序列
\SourceDiagram{diagram-3-15}{0.61061473}
\end{corollary}
反过来不成立，即上述推论中的性质不足以刻画正合三角（利用\stref{笔记}{3.12}，很容易构造反例）。\znote{40}{由殊三角取 Hom 得到的两条长正合列，是三角结构最主要的推论：一条关于协变 Hom，一条关于反变 Hom，与拓扑中纤维序列给出的同伦群长正合列完全平行。反过来，仅凭长正合列并不能判定一个三角是殊三角，这正是本段强调反向不成立的意思。[来源：nLab "triangulated category"，https:/\allowbreak /\allowbreak ncatlab.org/\allowbreak nlab/\allowbreak show/\allowbreak triangulated+\allowbreak category（2026-11-15 访问）]}
\begin{corollary}\label{st:3.16}
取两个正合三角
\SourceEquationDiagram{diagram-3-14}{0.38007652}{3.14}
以及映射 $[b_0]\in\Hom^0_{H(\A)}(Y_0,\widetilde Y_0)$、$[b_1]\in\Hom^0_{H(\A)}(Y_1,\widetilde Y_1)$，使其满足
\[
[\mu^2_{\A}(\tilde c_1,b_0)]=[\mu^2_{\A}(b_1,c_1)].
\]
那么存在一个（不唯一的）态射 $[b_2]$，将图补全为\eref{3.14}所示的交换图。此外，若 $[b_0],[b_1]$ 是同构，那么 $[b_2]$ 也是同构。
\end{corollary}

\SourcePage{49}
\begin{corollary}\label{st:3.17}
假设已有正合三角\eref{3.9}，且 $Y_0$ 存在平移 $SY_0$。那么旋转后的三角
\SourceDiagram{diagram-3-17}{0.23053822}
仍然正合。
\end{corollary}
在上一推论的陈述中，我们隐含地使用了\eref{3.3}，把 $[c_3]$ 视为到 $SY_0$ 的态射。对 $[c_1]$，类似地使用\eref{3.5}，其中非平凡的部分是应用平移函子 $S$，故采用记号 $[S(c_1)]$。

\begin{note}\label{st:3.18}
刚才的推论有一个简单而有用的结论。假设正合三角\eref{3.9}中的映射 $[c_1]$ 为零。由于抽象映射锥的同构类不依赖于上闭链的选择，\stref{笔记}{3.5}中的结论仍成立，即在 $H^0(\A)$ 中有 $Y_2\cong SY_0\oplus Y_1$。通过旋转正合三角，当其他任意一个映射为零时，也能得到类似结论。完整列表为
\[
\begin{aligned}
{}[c_1]=0&\quad\Longrightarrow\quad Y_2\cong SY_0\oplus Y_1,\\
[c_2]=0&\quad\Longrightarrow\quad Y_0\cong Y_1\oplus S^{-1}Y_2,\\
[c_3]=0&\quad\Longrightarrow\quad Y_1\cong Y_2\oplus Y_0.
\end{aligned}
\]
\end{note}
\begin{corollary}\label{st:3.19}
给定正合三角
\SourceDiagram{diagram-3-19a}{0.62515318}
还存在正合三角
\SourceDiagram{diagram-3-19b}{0.13084601}
\end{corollary}
习惯上把三角范畴的公理列为 (TR1)–(TR4)\bcite{61}，第 5 章。\stref{推论}{3.17}对应于 (TR2)；\stref{推论}{3.15}与\stref{推论}{3.16}是 (TR1)–(TR3) 的推论；\stref{推论}{3.19}则是八面体公理 (TR4) 的一个不完整版本。

\SourcePage{50}
\subsection{生成元}\label{sec:3j}
设 $\B$ 是三角化的 $\Ai$ 范畴，$\A\subset\B$ 是非空的全 $\Ai$ 子范畴。
\begin{definition*}
考虑满足如下条件的最小全子范畴 $\widetilde{\B}\subset\B$：它包含 $\A$，对同构封闭（即若 $X_0\cong X_1$ 在 $H^0(\B)$ 中成立，且 $X_0\in\Ob\widetilde{\B}$，则 $X_1\in\Ob\widetilde{\B}$），且自身三角化。称它为 $\B$ 中\emph{由 $\A$ 生成的三角化子范畴}；等价地，也可以说由 $\Ob\A\subset\Ob\B$ 生成。若 $\widetilde{\B}=\B$，则称 $\A$ \emph{生成} $\B$，或称其对象为 $\B$ 的\emph{生成元}。
\end{definition*}
具体地，从 $\Ob\A$ 出发，构造所有可能的映射锥及平移（向上与向下），再把这个过程重复任意有限多次，就得到 $\Ob\widetilde{\B}$。这样得到的子范畴是三角化的，因为作为\stref{推论}{3.9}与\stref{推论}{3.10}的特殊情形，一个三角在 $H(\widetilde{\B})$ 中正合，当且仅当它在 $H(\B)$ 中正合。

\begin{definition*}
设 $\A$ 是非空 $\Ai$ 范畴。一个\emph{三角化包络}是二元组 $(\B,\F)$，其中 $\B$ 为三角化的 $\Ai$ 范畴，$\F:\A\to\B$ 为上同调全忠实函子，且 $\F$ 的像中的对象生成 $\B$。
\end{definition*}
\begin{lemma}\label{st:3.20}
三角化包络总是存在的。此外，若 $(\B,\F)$ 与 $(\widetilde{\B},\widetilde{\F})$ 是 $\A$ 的两个这样的包络，则存在拟等价 $\G:\B\to\widetilde{\B}$，使 $\G\circ\F$ 在 $H^0(\operatorname{fun}(\A,\widetilde{\B}))$ 中同构于 $\widetilde{\F}$。
\end{lemma}
特别地，三角范畴 $H^0(\B)$ 在等价意义下不依赖于包络的选择；事实上，由\stref{命题}{3.14}的后一部分，这是正合等价。
\begin{definition*}
称这个范畴为 $\A$ 的\emph{导出范畴}，记作 $D(\A)$。
\end{definition*}
这个定义的优点，是允许采用不同的三角化包络给出不同实现。不过，对大多数应用而言，这种一般性带来的妨碍多于帮助。因此，下面不再保持中立，而选定一种具体的包络，即扭曲复形的 $\Ai$ 范畴 $\Tw\A$。它的定义仿照加性范畴上的链复形所组成的 dg 范畴。在这个范畴中，映射锥运算有显式实现，因而容易研究；我们将以此为出发点，证明上述所有抽象性质。
\begin{note}\label{st:3.21}
至少还有一种典范的三角化包络构造，它基于\stref{例题}{3.13}：在 $\Q=\modcat(\A)$ 中，取由 Yoneda 嵌入的像生成的全子范畴即可。
\end{note}
\begin{note}\label{st:3.22}
这里对三角结构采用的“链层次”或“内在”观点，最早出现于\bcite{28}。该文还引入了扭曲复形，至少讨论了 dg 范畴的情形；推广到 $\Ai$ 范畴的工作则见于\bcite{90}。
\end{note}

\subsection{加性扩张}\label{sec:3k}
构造 $\Tw\A$ 的第一步是加性扩张 $\Sigma\A$。设 $\A$ 是非含幺 $\Ai$ 范畴。
\begin{definition*}
$\Sigma\A$ 的对象是三元组
\SourcePage{51}
$X=(I,\{X^i\},\{V^i\})$，更直观地写作形式直和
\begin{equation}\label{eq:3.15}
X=\bigoplus_{i\in I}V^i\otimes X^i,
\end{equation}
其中 $I$ 是有限集，$\{X^i\}_{i\in I}$ 是 $\A$ 的一族对象，$\{V^i\}_{i\in I}$ 是相应的一族有限维分次向量空间，称为\emph{重数空间}。两个这样的直和之间的态射，由其组成直和项之间的态射适当地张量相应向量空间之间的线性映射构成：
\begin{equation}\label{eq:3.16}
\begin{aligned}
\hom_{\Sigma\A}\left(\bigoplus_{i\in I_0}V_0^i\otimes X_0^i,\,\bigoplus_{j\in I_1}V_1^j\otimes X_1^j\right)
&=\bigoplus_{ij}\hom_{\K}(V_0^i,V_1^j)\otimes\hom_{\A}(X_0^i,X_1^j),
\end{aligned}
\end{equation}
并带有自然的分次。
\end{definition*}
把态射写作矩阵 $a=(a^{ji})$，行列指标分别为 $j\in I_1$、$i\in I_0$；每个矩阵元形如
\[
a^{ji}=\sum_k\phi^{jik}\otimes x^{jik}\in\hom_{\K}(V_0^i,V_1^j)\otimes\hom_{\A}(X_0^i,X_1^j).
\]
写公式时，常只考虑求和仅含一项的情形 $a^{ji}=\phi^{ji}\otimes x^{ji}$，并默认按线性扩张。有了这一约定，即可定义 $\Sigma\A$ 上的复合映射。它把 $\A$ 的复合与线性映射的通常复合结合起来，正如\stref{笔记}{1.11}中的张量积：
\begin{equation}\label{eq:3.17}
\begin{aligned}
\mu^d_{\Sigma\A}(a_d,\ldots,a_1)^{i_d,i_0}
&=\sum_{i_1,\ldots,i_{d-1}}(-1)^{\triangleleft}\phi_d^{i_d,i_{d-1}}\cdots\phi_1^{i_1,i_0}\\[-2pt]
&\qquad\otimes\mu^d_{\A}(x_d^{i_d,i_{d-1}},\ldots,x_1^{i_1,i_0}),
\end{aligned}
\end{equation}
其中 $\triangleleft=\sum_{p<q}|\phi_p^{i_p,i_{p-1}}|\cdot(|x_q^{i_q,i_{q-1}}|-1)$，$a_k=(a_k^{ji})=(\phi_k^{ji}\otimes x_k^{ji})$。从现在起固定一个单元素集 $\{*\}$，把 $\A$ 的对象等同于三元组 $(\{*\},\{X^*=X\},\{V^*=\K\})$，其中 $\K$ 集中在零次。这样就把 $\A$ 嵌入为 $\Sigma\A$ 的全 $\Ai$ 子范畴。

\subsection{扭曲复形}\label{sec:3l}
\begin{definition*}
$\A$ 中的一个\emph{预扭曲复形}是二元组 $(X,\delta_X)$，其中 $X\in\Ob\Sigma\A$，而微分（有时也称联络）为 $\delta_X\in\hom^1_{\Sigma\A}(X,X)$。展开来写，就是
\begin{equation}\label{eq:3.18}
\left(X=\bigoplus_{i\in I}V^i\otimes X^i,\quad\delta_X=(\delta_X^{ji})\right),
\end{equation}
其中 $\delta_X^{ji}=\sum_k\phi^{jik}\otimes x^{jik}$，$\phi^{jik}\in\Hom_{\K}(V^i,V^j)$、$x^{jik}\in\hom_{\A}(X^i,X^j)$，且 $|\phi^{jik}|+|x^{jik}|=1$。
\end{definition*}
出于技术上的原因，还需要引入（简单的）子复形 $\widetilde X\subset X$ 与商复形 $X/\widetilde X$ 的概念。
\begin{definition*}
子复形的定义是：对每个 $i\in I$，取分次子空间 $\widetilde V^i\subset V^i$，使所有 $\phi^{jik}$ 都把 $\widetilde V^i$ 映入 $\widetilde V^j$，然后取 $\bigoplus_i\widetilde V^i\otimes X^i$，配以相应限制的微分。
\SourcePage{52}
商复形以相同方式构造，使用向量空间 $V^i/\widetilde V^i$ 以及 $\phi^{jik}$ 在它们之间诱导的线性映射。
\end{definition*}
\begin{definition*}
\emph{扭曲复形}是具有如下两个性质的预扭曲复形 $(X,\delta_X)$。首先，$\delta_X$ \emph{严格下三角}，即存在有限递减的子复形滤过
\[
X=F^0X\supset F^1X\supset\cdots\supset F^nX=0,
\]
使商 $F^kX/F^{k+1}X$ 上的诱导微分为零；滤过并不是 $(X,\delta_X)$ 结构的一部分，只要求它存在。其次，满足\emph{广义 Maurer–Cartan 方程}
\begin{equation}\label{eq:3.19}
\sum_{r=1}^{\infty}\mu^r_{\Sigma\A}(\delta_X,\ldots,\delta_X)=0.
\end{equation}

\end{definition*}
这个式子有意义，因为结合下三角条件与定义\eref{3.17}可见，\eref{3.19}中除有限多项外，其余各项均为零。常常省略记号中的微分，直接写作 $X\in\Ob\Tw\A$。

扭曲复形组成非含幺 $\Ai$ 范畴 $\Tw\A$ 的对象。态射空间与之前相同，即 $\hom_{\Tw\A}(X_0,X_1)=\hom_{\Sigma\A}(X_0,X_1)$，但所有复合都由微分的贡献作了变形：
\begin{equation}\label{eq:3.20}
\begin{aligned}
\mu^d_{\Tw\A}(a_d,\ldots,a_1)
&=\sum_{i_0,\ldots,i_d}\mu^{d+i_0+\cdots+i_d}_{\Sigma\A}\bigl(\overbrace{\delta_{X_d},\ldots,\delta_{X_d}}^{i_d},a_d,\\[-2pt]
&\qquad\underbrace{\delta_{X_{d-1}},\ldots,\delta_{X_{d-1}}}_{i_{d-1}},a_{d-1},\ldots,a_1,\underbrace{\delta_{X_0},\ldots,\delta_{X_0}}_{i_0}\bigr),
\end{aligned}
\end{equation}
其中对所有 $i_0,\ldots,i_d\geq0$ 求和。同样，微分的下三角性保证，插入足够多个 $\delta_{X_k}$ 后，复合就为零，因此\eref{3.20}是有限和。$\Tw\A$ 中结合律等式\eref{1.2}的成立，恰好归结为广义 Maurer–Cartan 方程\eref{3.19}。还要注意，微分为零的扭曲复形所组成的全子范畴可以与 $\Sigma\A$ 等同。

设 $X_0,X_1$ 是配有上述滤过的扭曲复形。这些滤过诱导链复形 $\hom_{\Tw\A}(X_0,X_1)$ 的滤过，从而给出收敛到其上同调的谱序列，起始项为
\begin{equation}\label{eq:3.21}
E_1^{rs}=\bigoplus_{j-i=r}\Hom^{r+s}_{H(\Sigma\A)}(F^iX_0/F^{i+1}X_0,\,F^jX_1/F^{j+1}X_1).
\end{equation}
该谱序列与 $H(\Tw\A)$ 中的复合相容。这意味着，对任意三个扭曲复形 $X_0,X_1,X_2$，映射 $\mu^2_{\Tw\A}$ 保持滤过，从而在相应谱序列上诱导乘积。

\SourcePage{53}
\subsection{函子性}\label{sec:3m}
设 $\G:\A\to\B$ 为非含幺 $\Ai$ 函子。诱导函子 $\Sigma\G:\Sigma\A\to\Sigma\B$ 定义为
\begin{equation}\label{eq:3.22}
\begin{aligned}
(\Sigma\G)\left(\bigoplus_{i\in I}V^i\otimes X^i\right)&=\bigoplus_{i\in I}V^i\otimes\G(X^i),\\
(\Sigma\G)^d(a_d,\ldots,a_1)^{i_d,i_0}
&=\sum_{i_1,\ldots,i_{d-1}}(-1)^{\triangleleft}\phi_d^{i_d,i_{d-1}}\cdots\phi_1^{i_1,i_0}\\[-2pt]
&\qquad\otimes\G^d(x_d^{i_d,i_{d-1}},\ldots,x_1^{i_1,i_0}),
\end{aligned}
\end{equation}
其中 $\triangleleft$ 与 $a_k$ 同\eref{3.17}。此外，任意 $T\in\hom_{\nufun(\A,\B)}(\G_0,\G_1)$ 都诱导 $\Sigma^1T\in\hom_{\nufun(\Sigma\A,\Sigma\B)}(\Sigma\G_0,\Sigma\G_1)$，同样由下式给出：
\begin{equation}\label{eq:3.23}
\begin{aligned}
(\Sigma^1T)^d(a_d,\ldots,a_1)^{i_d,i_0}
&=\sum_{i_1,\ldots,i_{d-1}}(-1)^{\triangleleft}\phi_d^{i_d,i_{d-1}}\cdots\phi_1^{i_1,i_0}\\[-2pt]
&\qquad\otimes T^d(x_d^{i_d,i_{d-1}},\ldots,x_1^{i_1,i_0}).
\end{aligned}
\end{equation}
公式\eref{3.22}、\eref{3.23}连同为零的高阶项，定义一个非含幺 $\Ai$ 函子 $\Sigma:\nufun(\A,\B)\to\nufun(\Sigma\A,\Sigma\B)$。

扭曲复形的情形完全类似。有一个诱导的非含幺 $\Ai$ 函子 $\Tw\G:\Tw\A\to\Tw\B$，
\[
\begin{aligned}
(\Tw\G)(X,\delta_X)&=\left(\Sigma\G(X),\sum_{e\geq1}(\Sigma\G)^e(\delta_X,\ldots,\delta_X)\right),\\[2pt]
(\Tw\G)^d(a_d,\ldots,a_1)
&=\sum_{i_0,\ldots,i_d}(\Sigma\G)^{d+i_0+\cdots+i_d}\bigl(\overbrace{\delta_{X_d},\ldots,\delta_{X_d}}^{i_d},a_d,\\[-2pt]
&\qquad\underbrace{\delta_{X_{d-1}},\ldots,\delta_{X_{d-1}}}_{i_{d-1}},a_{d-1},\ldots,a_1,\underbrace{\delta_{X_0},\ldots,\delta_{X_0}}_{i_0}\bigr).
\end{aligned}
\]
注意，若 $X$ 有子复形滤过 $F^kX$，使微分严格下三角，那么 $\Sigma\G(X)$ 由 $\Sigma\G(F^kX)$ 诱导的滤过也有同样性质，故它确实是扭曲复形。对 $T\in\hom_{\nufun(\A,\B)}(\G_0,\G_1)$，定义 $\Tw^1T\in\hom_{\nufun(\Tw\A,\Tw\B)}(\Tw\G_0,\Tw\G_1)$ 如下：
\begin{equation}\label{eq:3.24}
\begin{aligned}
(\Tw^1T)^d(a_d,\ldots,a_1)
&=\sum_{i_0,\ldots,i_d}(\Sigma^1T)^{d+i_0+\cdots+i_d}\bigl(\overbrace{\delta_{X_d},\ldots,\delta_{X_d}}^{i_d},a_d,\\[-2pt]
&\qquad\underbrace{\delta_{X_{d-1}},\ldots,\delta_{X_{d-1}}}_{i_{d-1}},a_{d-1},\ldots,a_1,\underbrace{\delta_{X_0},\ldots,\delta_{X_0}}_{i_0}\bigr).
\end{aligned}
\end{equation}
与之前相同，这些数据合在一起构成非含幺 $\Ai$ 函子 $\Tw:\nufun(\A,\B)\to\nufun(\Tw\A,\Tw\B)$，并且与左、右复合相容。这意味着，对固定的 $\G:\A\to\B$，有
\begin{equation}\label{eq:3.25}
\begin{aligned}
\LG_{\Tw\G}\circ\Tw&=\Tw\circ\LG_{\G}:\nufun(\C,\A)\longrightarrow\nufun(\Tw\C,\Tw\B),\\
\RG_{\Tw\G}\circ\Tw&=\Tw\circ\RG_{\G}:\nufun(\B,\C)\longrightarrow\nufun(\Tw\A,\Tw\C).
\end{aligned}
\end{equation}

\SourcePage{54}
\begin{lemma}\label{st:3.23}
若 $\G:\A\to\B$ 上同调全忠实，那么 $\Tw\G:\Tw\A\to\Tw\B$ 也上同调全忠实。
\end{lemma}
\begin{proof}
对 $(\Tw\G)^1$ 在谱序列\eref{3.21}上的作用使用比较论证，就把结论归结为 $\Sigma\G$ 的相应结论，而后者显然成立。
\end{proof}

\subsection{单位元与拟等价}\label{sec:3n}
若 $\A$ 是严格含幺 $\Ai$ 范畴，那么 $\Tw\A$ 也严格含幺。具体地，扭曲复形\eref{3.18}的严格恒等态射是“单位矩阵” $(\phi^{ji}\otimes x^{ji})$，其中 $\phi^{ii}\otimes x^{ii}=\id_{V^i}\otimes e_{X^i}$，而所有非对角矩阵元均为零。此外，若 $\G:\A\to\B$ 是严格含幺 $\Ai$ 函子，那么 $\Tw\G$ 也严格含幺。最后，对严格含幺的 $\A$ 及任意 $\C$，
\begin{equation}\label{eq:3.26}
\Tw:\nufun(\C,\A)\longrightarrow\nufun(\Tw\C,\Tw\A)
\end{equation}
是严格含幺 $\Ai$ 函子。只需回顾\secref{2b}中 $E_{\G}$ 的表达式，并代入\eref{3.24}，即可看出这一点。其上同调含幺版本则不那么显然：
\begin{lemma}\label{st:3.24}
\textup{(i)} 若 $\A$ 为 c-含幺，那么 $\Tw\A$ 也为 c-含幺。

\textup{(ii)} 假设 $\A,\B$ 都为 c-含幺，且 $\G:\A\to\B$ 是 c-含幺 $\Ai$ 函子。那么 $\Tw\G$ 也为 c-含幺。

\textup{(iii)} 对 c-含幺的 $\A$ 及任意 $\C$，\eref{3.26}是 c-含幺 $\Ai$ 函子。
\end{lemma}
\begin{proof}
(i) 如\stref{引理}{2.3}的证明，先取最小模型，再将其严格化，得到严格含幺的 $\widetilde{\A}$ 及在对象上为恒等映射的拟同构 $\Phi:\A\to\widetilde{\A}$。由\stref{引理}{3.23}，诱导的非含幺 $\Ai$ 函子 $\Tw\Phi:\Tw\A\to\Tw\widetilde{\A}$ 上同调全忠实。后者严格含幺；其恒等态射可经由这一全忠实函子在上同调上唯一拉回，故前者必然上同调含幺。

(ii) 固定 $X\in\Tw\A$ 及适当的滤过 $F^kX$。在保滤过的自同态子复形上，$(\Tw\Phi)^1$ 仍是拟同构；这由\eref{3.21}的同一个有限滤过比较论证得出。因此可以选择代表 $X$ 的上同调恒等态射的保滤过上闭链 $e_X$，使其在各相继商上诱导上同调恒等态射。由定义，$(\Tw\G)^1(e_X)$ 对 $\Tw\G(F^kX)$ 给出的像滤过是下三角的，且在商 $\Tw\G(F^kX)/\Tw\G(F^{k+1}X)$ 上诱导的自同态是恒等态射的链代表，因为 $\G$ 上同调含幺。利用\eref{3.21}可知，以 $[(\Tw\G)^1(e_X)]$ 作左乘和右乘，在 $\Hom_{H(\Tw\B)}(\Tw\G(X),\Tw\G(X))$ 上都是自同构。因此该元素本身可逆。它还是幂等元，因为它是恒等态射在一个函子下的像，所以它等于恒等态射。这证明 $\Tw\G$ 上同调含幺。

\SourcePage{55}
(iii) 固定非含幺 $\Ai$ 函子 $\G:\C\to\A$，令 $\widetilde{\G}=\Phi\circ\G$，其中 $\Phi$ 如上。作为\eref{3.25}的特殊情形，有交换图
\[
\begin{array}{ccc}
\hom_{\nufun(\C,\A)}(\G,\G)&\xrightarrow{\ \LG^1_{\Phi}\ }&\hom_{\nufun(\C,\widetilde{\A})}(\widetilde{\G},\widetilde{\G})\\[3pt]
\Big\downarrow\scriptstyle\Tw^1&&\Big\downarrow\scriptstyle\Tw^1\\[3pt]
\hom_{\nufun(\Tw\C,\Tw\A)}(\Tw\G,\Tw\G)&\xrightarrow{\ \LG^1_{\Tw\Phi}\ }&\hom_{\nufun(\Tw\C,\Tw\widetilde{\A})}(\Tw\widetilde{\G},\Tw\widetilde{\G}).
\end{array}
\]
由 (ii) 及\stref{引理}{2.4}，水平箭头在上同调层次上都把恒等态射映到恒等态射。右侧竖直箭头处有严格单位元，故它也如此。下方水平箭头由\stref{引理}{3.23}与\stref{引理}{1.7}在上同调上为同构，因此交换性给出所需结论。
\end{proof}

\begin{lemma}\label{st:3.25}
设 $\G:\A\to\B$ 是 c-含幺 $\Ai$ 范畴之间的拟等价。那么 $\Tw\G:\Tw\A\to\Tw\B$ 也是拟等价。
\end{lemma}
这是\stref{定理}{2.9}及上一引理的直接推论：先在相应函子范畴中的同构意义下取 $\G$ 的逆，再考察诱导的函子及自然变换即可。另一种较少依赖形式体系、更多借助正合三角的方法，见后文\stref{引理}{3.36}的证明。

\subsection{基本运算}\label{sec:3o}
从这里起，恢复之前的通常假设：所涉及的一切都上同调含幺。

$\Sigma\A$ 对象的形式允许定义直和运算 $Y_0\oplus Y_1$：形式上，把指标集及其相应的族作不交并即可。这个运算显然延伸到扭曲复形。此外，$Y_0\oplus Y_1$ 在 Yoneda 嵌入下的像，恰好就是\secref{3b}中的 $\Ai$ 模直和 $\Y_0\oplus\Y_1$。

同样自然地，可以对有限维分次向量空间 $Z$ 与 $\Sigma\A$ 中对象 $Y=\bigoplus_{i\in I}W^i\otimes Y^i$ 取张量积 $Z\otimes Y$。它有相同的指标集 $I$ 与对象 $Y^i$，但将每个 $W^i$ 替换为 $Z\otimes W^i$。给定 $a=(a^{ji})=(\phi^{ji}\otimes x^{ji})\in\hom_{\Sigma\A}(Y_0,Y_1)$，把每个线性映射 $\phi^{ji}\in\Hom_{\K}(W_0^i,W_1^j)$ 替换为分次张量积 $\id_Z\otimes\phi^{ji}\in\Hom_{\K}(Z\otimes W_0^i,Z\otimes W_1^j)$，即可定义 $\id_Z\otimes a\in\hom_{\Sigma\A}(Z\otimes Y_0,Z\otimes Y_1)$；这里
\[
(\id_Z\otimes\phi^{ji})(z\otimes w)=(-1)^{|\phi^{ji}|\cdot|z|}z\otimes\phi^{ji}(w).
\]
这给出一个高阶项为零的非含幺 $\Ai$ 函子 $Z\otimes-:\Sigma\A\to\Sigma\A$。扭曲复形也有类似构造，它把 $(Y,\delta_Y)$ 映到 $(Z\otimes Y,\delta_{Z\otimes Y}=\id_Z\otimes\delta_Y)$，并以与之前相同的方式作用于态射。在 Yoneda 嵌入下，它成为\secref{3c}中的 $\modcat(\Tw\A)$ 上的函子 $Z\otimes-$，这里限制在 $Z$ 有限维且 $\partial_Z=0$ 的情形。

\SourcePage{56}
当 $Z=\K[\sigma]$ 时，记 $S^{\sigma} Y=Z\otimes Y$、$S^{\sigma}(a)=\id_Z\otimes a$。把扭曲复形上的运算显式写出，也许会有所帮助：
\[
\begin{gathered}
\left(Y=\bigoplus_iW^i\otimes Y^i,\quad\delta_Y=\left(\sum_k\phi^{jik}\otimes x^{jik}\right)_{ji}\right)\\
\Big\downarrow\\
\left(S^{\sigma} Y=\bigoplus_iW^i[\sigma]\otimes Y^i,\quad
\delta_{S^{\sigma} Y}=S^{\sigma}(\delta_Y)=\left(\sum_k(-1)^{\sigma|\phi^{jik}|}\phi^{jik}\otimes x^{jik}\right)_{ji}\right).
\end{gathered}
\]
将前面的讨论专门用于此情形，即可看出：$S^{\sigma}$ 作为阶数 $\geq2$ 的项均为零的 $\Ai$ 函子，确实表示\secref{3d}中的抽象平移函子 $\ModShift^{\sigma}$。事实上，在链层次上有自然的等同
\[
\hom_{\Tw\A}(Y_0,SY_1)\cong\hom_{\Tw\A}(Y_0,Y_1)[1],
\]
它诱导\eref{3.3}。

作为补充，现在简要考察 $\A$ 严格含幺的情形。此时可以如下定义有限维向量空间链复形 $(Z,\partial_Z)$ 与扭曲复形 $Y=(\bigoplus_{i\in I}W^i\otimes Y^i,\delta_Y)$ 的张量积 $Z\otimes Y$。首先把微分改为 $\tilde\partial_Z(z)=(-1)^{|z|-1}\partial_Z(z)$，然后令
\begin{equation}\label{eq:3.27}
Z\otimes Y=\left(\bigoplus_{i\in I}(Z\otimes W^i)\otimes Y^i,\quad\id_Z\otimes\delta_Y+\tilde\partial_Z\otimes e_Y\right),
\end{equation}
其中微分的第二项定义为对角矩阵，其矩阵元为 $(\tilde\partial_Z\otimes\id_{W^i})\otimes e_{Y^i}$。直接计算可验证广义 Maurer–Cartan 方程；这里关键地用到了 $e_{Y^i}$ 是严格单位元。显然，$\partial_Z$ 相对于 $Z$ 的按次数滤过严格下三角。将这个滤过与 $Y$ 的滤过结合，就得到\eref{3.27}的滤过，从而完成验证：所得对象确实是扭曲复形。

\subsection{映射锥}\label{sec:3p}
设 $c\in\hom_{\Tw\A}(Y_0,Y_1)$ 是零次上闭链。
\begin{definition*}
映射锥 $C=\Cone(c)$ 是扭曲复形
\begin{equation}\label{eq:3.28}
\left(C=SY_0\oplus Y_1,\quad\delta_C=\begin{pmatrix}S(\delta_{Y_0})&0\\-S(c)&\delta_{Y_1}\end{pmatrix}\right).
\end{equation}
\end{definition*}
这里先用平移函子把 $c$ 变为元素 $S(c)\in\hom^0_{\Tw\A}(SY_0,SY_1)$，再以直接的方式把该空间与 $\hom^1_{\Tw\A}(SY_0,Y_1)$ 等同；比较\eref{3.5}。若每个 $Y_k$ 都有子复形滤过 $F^iY_k$，使微分严格下三角，那么映射锥的组合滤过——先取 $F^iSY_0\oplus Y_1$，再接上 $0\oplus F^iY_1$——对 $C$ 也具有同样的性质。广义 Maurer–Cartan 方程等价于各 $Y_k$ 的对应方程，以及 $c$ 是 $\Tw\A$ 中上闭链这一事实。$C$ 在 Yoneda 嵌入下的像可以与抽象映射锥 $\C=\ACone(c)$ 等同，只需利用
\[
\begin{aligned}
\hom_{\Tw\A}(X,C)&=\hom_{\Tw\A}(X,SY_0)\oplus\hom_{\Tw\A}(X,Y_1)\\
&\cong\hom_{\Tw\A}(X,Y_0)[1]\oplus\hom_{\Tw\A}(X,Y_1)=\C(X).
\end{aligned}
\]
\SourcePage{57}
综上，我们已经看到，$\Tw\A$ 含有表示下列运算的优选对象：直和；与有限维分次向量空间的张量积，特别是平移；在严格含幺情形下，还包括与有限维链复形的张量积；以及最重要的映射锥。以后处理扭曲复形时，自然都会采用这些代表。\znote{49}{这些优选代表使扭曲复形范畴成为“链层面可计算”的三角范畴：直和、张量与平移都有具体模型，映射锥则由\eref{3.28}给出。值得注意的是，在一般的三角范畴中映射锥并非函子性构造，本书用扭曲复形把它具体化，正是为了避免这一缺陷。[来源：nLab "triangulated category"，https:/\allowbreak /\allowbreak ncatlab.org/\allowbreak nlab/\allowbreak show/\allowbreak triangulated+\allowbreak category（2026-11-15 访问）]}

\begin{note}\label{st:3.26}
还可以换一种方式定义加性扩张：只允许一维重数空间，更具体地说，只允许 $\K$ 的平移。此时对象 $X$ 写作直和 $\bigoplus_iS^{-\sigma^i}X^i$，其中 $X^i\in\Ob\A$，$\sigma^i\in\Z$。用显然的同构，并对每个系数 $a^{ji}$ 加上额外的符号 $(-1)^{(|a^{ji}|-1)\sigma_0^i}$，来作如下等同：
\begin{equation}\label{eq:3.29}
\hom_{\Sigma\A}\left(\bigoplus_iS^{-\sigma_0^i}X_0^i,\,\bigoplus_jS^{-\sigma_1^j}X_1^j\right)
=\bigoplus_{i,j}\hom_{\A}(X_0^i,X_1^j)[\sigma_1^j-\sigma_0^i].
\end{equation}
这里 $|a^{ji}|$ 指施加任何平移之前，在 $\hom_{\A}(X_0^i,X_1^j)$ 中的次数。相对于这一等同，复合变为
\[
\mu^d_{\Sigma\A}(a_d,\ldots,a_1)^{i_d,i_0}
=(-1)^{\sigma_0^{i_0}}\sum_{i_1,\ldots,i_{d-1}}\mu^d_{\A}(a_d^{i_d,i_{d-1}},\ldots,a_1^{i_1,i_0}),
\]
其中 $a_k\in\hom_{\Sigma\A}(X_{k-1},X_k)$，$\sigma_0^{i_0}$ 是 $X_0$ 的第 $i_0$ 个直和项的平移量。平移函子 $S$ 将态射 $a\in\hom^k_{\Sigma\A}(X_0,X_1)$ 乘以 $(-1)^{k-1}$。若沿同样思路继续定义一个版本的 $\Tw\A$，则 $C=\Cone(c)$ 上的微分为
\[
\delta_C=\begin{pmatrix}\delta_{Y_0}&0\\c&\delta_{Y_1}\end{pmatrix};
\]
\eref{3.28}中的平移运算与符号已经吸收到等同\eref{3.29}中。在拟等价意义下，允许一般的重数空间，还是只允许 $\K$ 的平移，当然并无区别。较一般的方法，其优点是无需选择基就能进行张量积运算 $Z\otimes-$，而且符号也更系统，尽管稍微复杂一些；不过，总体差别很小。
\end{note}

